\section{Limitations and further work}

\begin{description}[leftmargin=2.2em,style=nextline]
\item[Single field age.] $\Rg$ is evaluated at one $t_{\mathrm{op}}$. The
decay of the loss from the first year to the thirtieth is obtained by repeated
evaluation rather than by an outer loop over years, and the warming of the
rock within the cluster is not represented as a capacitance. At island spacing
that warming is a one-time charge of order \SI{360}{\mega\watt\hour} per well,
roughly 300 daily cycles, after which it is storage rather than loss.
\item[The near-well resistance.] $R_p$ ignores the crowding of the conduction
paths near the casing. It is a quarter of a total the ground dominates.
\item[Uniform cluster.] Hypothesis B6 divides the perimeter loss equally. Edge
wells lose more and interior wells less; representing that requires a
position-resolved field rather than a single $R_f(z)$.
\item[The surface machinery is frozen.] Round-trip efficiency is governed by
the heat-pump lift. Writing the two Carnot factors together,
\begin{equation}
\mathrm{COP}_C \times \eta_C
 = \frac{T_h}{T_h - T_{\mathrm{src}}}\left(1 - \frac{T_0}{T_h}\right)
 = \frac{T_h - T_0}{T_h - T_{\mathrm{src}}},
\end{equation}
which depends on the \emph{lift} and not on how hot the store is --- so
raising $\Tm$ reduces $\etarte$, because the coefficient of performance falls
faster than the power block gains. The lever is the source temperature, and
the same field is considerably hotter at depth than the \SI{60}{\degC}
currently assumed. That is a surface-side study and is not addressed here.
\end{description}
